% ---------------------------------------------------------------------------
%  Diagrama entidad-relación, con la frontera de cada base marcada.
%
%  Un solo DER global y no uno por base, a propósito: lo que se evalúa en la
%  sección 5 es la independencia de datos, y eso solo se ve cuando las tres
%  bases aparecen juntas y se distingue qué referencias cruzan la frontera y
%  cuáles no.
%
%  Compilar:  pdflatex -output-directory=../figuras modelo-datos.tex
%  (el PDF resultante es informe/figuras/modelo-datos.pdf)
% ---------------------------------------------------------------------------
\documentclass[border=6pt]{standalone}
\usepackage[T1]{fontenc}
\usepackage[utf8]{inputenc}
\usepackage[spanish]{babel}
\usepackage{lmodern}
\usepackage{tikz}
\usetikzlibrary{positioning, fit, backgrounds, calc}

\definecolor{bordeBase}{HTML}{9AA3AF}
\definecolor{fondoBase}{HTML}{F3F4F6}
\definecolor{cabecera}{HTML}{1F3A5F}
\definecolor{tablaFondo}{HTML}{FFFFFF}
\definecolor{acento}{HTML}{B45309}
\definecolor{textoTenue}{HTML}{4B5563}

\begin{document}
\begin{tikzpicture}[
    font=\sffamily\footnotesize,
    tabla/.style={
      rectangle, draw=cabecera!60, thick, fill=tablaFondo,
      inner sep=0pt, align=left, text width=52mm, rounded corners=1.5pt},
    titulo/.style={font=\sffamily\bfseries\small, text=white},
    fk/.style={-latex, draw=cabecera, thick},
    logica/.style={-latex, draw=acento, thick, dashed},
    nota/.style={font=\sffamily\scriptsize\itshape, text=textoTenue, align=center},
  ]

  % --- Macro para dibujar una tabla: cabecera de color + lista de columnas ---
  \newcommand{\entidad}[4]{% #1 nodo, #2 posición, #3 nombre, #4 columnas
    \node[tabla] (#1) #2 {%
      \begin{minipage}{52mm}
        \colorbox{cabecera}{\makebox[52mm][l]{\hspace{2mm}\titulostyle{#3}}}\\[1mm]
        \hspace*{2mm}\begin{minipage}{48mm}#4\end{minipage}\\[1mm]
      \end{minipage}};}
  \newcommand{\titulostyle}[1]{\textcolor{white}{\sffamily\bfseries\small #1}}
  \newcommand{\pk}[1]{\textbf{#1}\,\textcolor{cabecera}{\scriptsize PK}}
  \newcommand{\uk}[1]{#1\,\textcolor{cabecera}{\scriptsize UQ}}
  \newcommand{\rf}[1]{#1\,\textcolor{acento}{\scriptsize ref}}

  % ======================= usuarios_db =======================
  \entidad{usuario}{at (0,0)}{usuario}{%
    \pk{id} : bigint\\
    \uk{username} : varchar(60)\\
    nombre : varchar(120)\\
    password\_hash : varchar(72)\\
    rol : varchar(20) \textsc{check}\\
    activo : boolean\\
    creado\_en : timestamp}

  % ======================= canchas_db ========================
  \entidad{cancha}{at (7.4,0)}{cancha}{%
    \pk{id} : bigint\\
    nombre : varchar(120)\\
    deporte : varchar(20) \textsc{check}\\
    hora\_apertura : time\\
    hora\_cierre : time \textsc{check}\\
    activa : boolean}

  \entidad{bloqueo}{at (7.4,-4.6)}{bloqueo\_mantenimiento}{%
    \pk{id} : bigint\\
    cancha\_id : bigint \textcolor{cabecera}{\scriptsize FK}\\
    desde : timestamp\\
    hasta : timestamp \textsc{check}\\
    motivo : varchar(200)}

  % ======================= reservas_db =======================
  \entidad{reserva}{at (14.8,0)}{reserva}{%
    \pk{id} : bigint\\
    \rf{cancha\_id} : bigint\\
    \rf{usuario\_id} : bigint\\
    fecha : date\\
    hora\_inicio : time\\
    hora\_fin : time \textsc{check}\\
    estado : varchar(20) \textsc{check}\\
    creada\_en : timestamp\\
    cancelada\_en : timestamp}

  \entidad{configuracion}{at (14.8,-5.6)}{configuracion}{%
    \pk{clave} : varchar(60)\\
    valor : varchar(120)}

  % --------- Fronteras de base de datos ---------
  \begin{scope}[on background layer]
    \node[draw=bordeBase, dashed, thick, fill=fondoBase, rounded corners=3pt,
          fit=(usuario), inner sep=7mm] (bd1) {};
    \node[draw=bordeBase, dashed, thick, fill=fondoBase, rounded corners=3pt,
          fit=(cancha)(bloqueo), inner sep=7mm] (bd2) {};
    \node[draw=bordeBase, dashed, thick, fill=fondoBase, rounded corners=3pt,
          fit=(reserva)(configuracion), inner sep=7mm] (bd3) {};
  \end{scope}

  \node[font=\sffamily\bfseries\small, text=textoTenue, above=1mm of bd1.north] {usuarios\_db · ms-usuarios};
  \node[font=\sffamily\bfseries\small, text=textoTenue, above=1mm of bd2.north] {canchas\_db · ms-canchas};
  \node[font=\sffamily\bfseries\small, text=textoTenue, above=1mm of bd3.north] {reservas\_db · ms-reservas};

  % --------- Relaciones ---------
  % Dentro de una base: clave foránea de verdad.
  \draw[fk] (bloqueo.north) -- node[right=1mm, font=\sffamily\scriptsize] {FK} (cancha.south);

  % Cruzando la frontera: referencia lógica, SIN clave foránea posible.
  \draw[logica] (reserva.west) to[out=180, in=0]
    node[above, font=\sffamily\scriptsize, pos=0.5] {cancha\_id}
    node[below, font=\sffamily\scriptsize, pos=0.5] {(sin FK)} (cancha.east);

  % usuario_id va por debajo de todo, a una altura fija: si se enrutara en
  % curva libre cruzaría la caja de bloqueos y taparía sus columnas.
  \coordinate (bajo) at (0, -11.2);
  \draw[logica] (bd3.south west) -- (bd3.south west |- bajo)
    -- node[above, font=\sffamily\scriptsize, pos=0.5]
         {usuario\_id · referencia lógica, sin clave foránea: cruza la frontera de base}
       (usuario.south |- bajo)
    -- (usuario.south);

  % --------- Nota de la restricción EXCLUDE ---------
  \node[nota, align=left, text width=62mm, below=5mm of bd3.south east, anchor=north east] (excl) {%
    \textcolor{acento}{\textbf{reserva\_sin\_solape}} · \textsc{exclude using gist}\\
    sobre (cancha\_id, tsrange(fecha + hora, '[)'))\\
    \textsc{where} estado = 'CONFIRMADA'.\\[0.6mm]
    Implementa RN-02 también ante peticiones concurrentes:\\
    es lo que ninguna comprobación previa puede garantizar.};
  \draw[acento, thick, dotted] (excl.north) -- (excl.north |- bd3.south);

  % --------- Leyenda ---------
  % Las muestras de línea se dibujan, no se simulan con \rule: así la de
  % referencia lógica se ve discontinua, que es justo lo que la distingue.
  \node[anchor=north west, font=\sffamily\scriptsize, text=textoTenue, align=left]
        at ($(bd1.west |- bajo) + (0,-1.2)$) {%
    \tikz[baseline=-0.5ex]{\draw[fk] (0,0) -- (0.7,0);}\; clave foránea, dentro de una misma base\\[1mm]
    \tikz[baseline=-0.5ex]{\draw[logica] (0,0) -- (0.7,0);}\; referencia lógica, cruza la frontera entre bases\\[2mm]
    \begin{minipage}{112mm}\itshape
      Las referencias que cruzan la frontera no pueden ser claves foráneas: apuntan a
      tablas de otra base, a la que esa conexión ni siquiera llega. Las valida
      \textsc{BookingService} llamando a \textsc{ms-canchas} antes de persistir, y por
      eso una reserva sobre una cancha inexistente se rechaza con 404 en lugar de
      quedar registrada.
    \end{minipage}};

\end{tikzpicture}
\end{document}
